\ifdefined\gkver\else\def\gkver{student}\fi
\documentclass[\gkver]{gaokaozhenti}
\begin{document}

\chapter{2013年山东理综}
%% paperType: 理综物理部分

\begin{enumerate}
\item
%% number: 14
%% paperName: 2013年山东理综第 14 题 5 分
%% typeId: 2
%% type: 多选题
%% chapter: 其他
%% point: 物理思想
%% method: 无
%% score: 5
%% degree: 850
%% duplicateId: 0
%% body:
伽利略开创了实验研究和逻辑推理相结合探索物理规律的科学方法，利用这种方法伽利略发现的规律有\xzanswer{AC}

\fourchoices[answer=AC]
{力不是维持物体运动的原因}
{物体之间普遍存在相互吸引力}
{忽略空气阻力，重物与轻物下落得同样快}
{物体间的相互作用力总是大小相等，方向相反}

%% answer: AC
%% memo:
\memoanswer{
本题原卷及材料中未提供详解，待补充。
}

\item
%% number: 15
%% paperName: 2013年山东理综第 15 题 5 分
%% typeId: 1
%% type: 单选题
%% chapter: 力物体的平衡
%% point: 物体系平衡
%% method: 无
%% score: 5
%% degree: 650
%% duplicateId: 0
%% body:
如图所示，用完全相同的轻弹簧A、B、C将两个相同的小球连接并悬挂，小球处于静止状态，弹簧A与竖直方向的夹角为$30\Udu$，弹簧C水平，则弹簧A、C的伸长量之比为\xzanswer{D}

\onepicture[w=4cm]{figs/15.png}

\fourchoices[answer=D]
{$\sqrt 3 :4$}
{$4:\sqrt 3$}
{$1:2$}
{$2:1$}

%% answer: D
%% memo:
\memoanswer{
将两小球及弹簧B视为一个整体系统，因小球处于静止状态，故该系统水平方向受力平衡，则有$k\Delta x_{A}\sin30\Udu=k\Delta x_{C}$，可得$\Delta x_{A}:\Delta x_{C}=2:1$，选项D正确。
}

\item
%% number: 16
%% paperName: 2013年山东理综第 16 题 5 分
%% typeId: 2
%% type: 多选题
%% chapter: 功和能
%% point: 动能定理
%% method: 无
%% score: 5
%% degree: 650
%% duplicateId: 0
%% body:
如图所示，楔形木块abc固定在水平面上，粗糙斜面ab和光滑斜面bc与水平面的夹角相同，顶角b处安装一定滑轮。质量分别为$M$、$m$（$M>m$）的滑块，通过不可伸长的轻绳跨过定滑轮连接，轻绳与斜面平行。两滑块由静止释放后，沿斜面做匀加速运动。若不计滑轮的质量和摩擦，在两滑块沿斜面运动的过程中\xzanswer{CD}

\onepicture[w=7cm]{figs/16.png}

\fourchoices[answer=CD]
{两滑块组成系统的机械能守恒}
{重力对$M$做的功等于$M$动能的增加}
{轻绳对$m$做的功等于$m$机械能的增加}
{两滑块组成系统的机械能损失等于$M$克服摩擦力做的功}

%% answer: CD
%% memo:
\memoanswer{
本题原卷及材料中未提供详解，待补充。
}

\item
%% number: 17
%% paperName: 2013年山东理综第 17 题 5 分
%% typeId: 2
%% type: 多选题
%% chapter: 交流电
%% point: 交流电
%% method: 无
%% score: 5
%% degree: 650
%% duplicateId: 0
%% body:
图甲是小型交流发电机的示意图，两磁极 N、S 间的磁场可视为水平方向的匀强磁场，A 为交流电流表。线圈绕垂直于磁场方向的水平轴 $OO'$ 沿逆时针方向匀速转动。从图示位置开始计时，产生的交变电流随时间变化的图像如图乙所示。以下判断正确的是\xzanswer{AC}

\onepicture[w=10cm]{figs/17.png}

\fourchoices[answer=AC]
{电流表的示数为$10\UA$}
{线圈转动的角速度为$50\pi\Urads$}
{0.01 s 时线圈平面与磁场方向平行}
{0.02 s 时电阻$R$中电流的方向自右向左}

%% answer: AC
%% memo:
\memoanswer{
本题原卷及材料中未提供详解，待补充。
}

\item
%% number: 18
%% paperName: 2013年山东理综第 18 题 5 分
%% typeId: 1
%% type: 单选题
%% chapter: 电磁感应
%% point: 电磁感应中图象
%% method: 无
%% score: 5
%% degree: 650
%% duplicateId: 0
%% body:
一段导线绕成图甲所示的闭合电路，并固定在水平面（纸面）内，回路的$ab$边置于垂直纸面向里的匀强磁场Ⅰ中。回路的圆形区域内有垂直纸面的磁场Ⅱ，以向里为磁场Ⅱ的正方向，其磁感应强度$B$随时间$t$变化的图像如图乙所示。用$F$表示$ab$边受到的安培力，以水平向右为$F$的正方向，能正确反映$F$随时间$t$变化的图像是\xzanswer{B}

\twopicture[numA=甲, numB=乙, hA=2.6cm, hB=2.6cm]
{figs/18a.png}
{figs/18b.png}

\fourchoices[answer=B, ispicture=true, h=2.2cm]
{figs/18c.png}
{figs/18d.png}
{figs/18e.png}
{figs/18f.png}

%% answer: B
%% memo:
\memoanswer{
本题原卷及材料中未提供详解，待补充。
}

\item
%% number: 19
%% paperName: 2013年山东理综第 19 题 5 分
%% typeId: 2
%% type: 多选题
%% chapter: 电场
%% point: 电场线
%% method: 无
%% score: 5
%% degree: 450
%% duplicateId: 0
%% body:
如图所示，在$x$轴上相距为$L$的两点固定两个等量异种点电荷$+Q$、$-Q$，虚线是以$+Q$所在点为圆心、$L/2$为半径的圆，$a$、$b$、$c$、$d$是圆上的四个点，其中$a$、$c$两点在$x$轴上，$b$、$d$两点关于$x$轴对称。下列判断正确的是\xzanswer{ABD}

\onepicture[w=5cm]{figs/19.png}

\fourchoices[answer=ABD]
{$b$、$d$两点处的电势相同}
{四点中$c$点处的电势最低}
{$b$、$d$两点处的电场强度相同}
{将一试探电荷$+q$沿圆周由$a$点移至$c$点，$+q$的电势能减小}

%% answer: ABD
%% memo:
\memoanswer{
本题原卷及材料中未提供详解，待补充。
}

\item
%% number: 20
%% paperName: 2013年山东理综第 20 题 5 分
%% typeId: 1
%% type: 单选题
%% chapter: 曲线运动
%% point: 天体运动
%% method: 无
%% score: 5
%% degree: 450
%% duplicateId: 0
%% body:
双星系统由两颗恒星组成，两恒星在相互引力的作用下，分别围绕其连线上的某一点做周期相同的匀速圆周运动。研究发现，双星系统演化过程中，两星的总质量、距离和周期均可能发生变化。若某双星系统中两星做圆周运动的周期为$T$，经过一段时间演化后，两星总质量变为原来的$k$倍，两星之间的距离变为原来的$n$倍，则运动的周期为\xzanswer{B}

\fourchoices[answer=B]
{$\sqrt{\frac{n^{3}}{k^{2}}}T$}
{$\sqrt{\frac{n^{3}}{k}}T$}
{$\sqrt{\frac{n^{2}}{k}}T$}
{$\sqrt{\frac{n}{k}}T$}

%% answer: B
%% memo:
\memoanswer{
本题原卷及材料中未提供详解，待补充。
}

\item
%% number: 21
%% paperName: 2013年山东理综第 21 题 9 分
%% typeId: 4
%% type: 实验题
%% chapter: 电路
%% point: 电路实验
%% method: 无
%% score: 9
%% degree: 450
%% duplicateId: 0
%% body:
（1）图甲为一游标卡尺的结构示意图，当测量一钢笔帽的内径时，应该用游标卡尺的\tkanswer[1.2cm]{A}（填“A”、“B”或“C”）进行测量；示数如图乙所示，该钢笔帽的内径为\tkanswer[1.6cm]{11.30}$\Umm$。

\onepicture[w=12cm]{figs/21.png}

（2）霍尔效应是电磁基本现象之一，近期我国科学家在该领域的实验研究上取得了突破性进展。如图丙所示，在一矩形半导体薄片的$P$、$Q$间通入电流$I$，同时外加与薄片垂直的磁场$B$，在$M$、$N$间出现电压$U_{H}$，这个现象称为霍尔效应，$U_{H}$称为霍尔电压，且满足$U_{H}=k\frac{IB}{d}$，式中$d$为薄片的厚度，$k$为霍尔系数。某同学通过实验来测定该半导体薄片的霍尔系数。

\onepicture[h=3.5cm, align=right]{figs/21a.png}

\begin{enumerate}
	\item
	若该半导体材料是空穴（可视为带正电粒子）导电，电流与磁场方向如图丙所示，该同学用电压表测量$U_{H}$时，应将电压表的“$+$”接线柱与\tkanswer[1cm]{$M$}（填“$M$”或“$N$”）端通过导线相连。

	\item
	已知薄片厚度$d=0.40\Umm$，该同学保持磁感应强度$B=0.10\UT$不变，改变电流$I$的大小，测量相应的$U_{H}$值，记录数据如下表所示。

	\begin{center}
	\begin{tabular}{|c|c|c|c|c|c|c|}
	\hline
	$I$（$\times10^{-3}\UA$） & 3.0 & 6.0 & 9.0 & 12.0 & 15.0 & 18.0 \\
	\hline
	$U_{H}$（$\times10^{-3}\UV$） & 1.1 & 1.9 & 3.4 & 4.5 & 6.2 & 6.8 \\
	\hline
	\end{tabular}
	\end{center}

	根据表中数据在给定区域内（见答题卡）画出$U_{H}-I$图线，利用图线求出该材料的霍尔系数为\tkanswer[2cm]{1.5}$\times10^{-3}\mathrm{V}\cdot\mathrm{m}\cdot\mathrm{A}^{-1}\cdot\mathrm{T}^{-1}$（保留2位有效数字）。

	\item
	该同学查阅资料发现，使半导体薄片中的电流反向再次测量，取两个方向测量的平均值，可以减小霍尔系数的测量误差，为此该同学设计了如图丁所示的测量电路，$S_{1}$、$S_{2}$均为单刀双掷开关，虚线框内为半导体薄片（未画出）。为使电流从$Q$端流入，$P$端流出，应将$S_{1}$掷向\tkanswer[1cm]{$b$}（填“$a$”或“$b$”），$S_{2}$掷向\tkanswer[1cm]{$c$}（填“$c$”或“$d$”）。为了保证测量安全，该同学改进了测量电路，将一合适的定值电阻串联在电路中。在保持其它连接不变的情况下，该定值电阻应串联在相邻器件\tkanswer[1cm]{$S_{1}$}和\tkanswer[1cm]{$E$}（填器件代号）之间。

	\onepicture[h=3.5cm, align=right]{figs/21b.png}
\end{enumerate}

%% answer:
\jdanswer{
\begin{enumerate}
	\item
	$M$

	\item
	1.5（1.4或1.6），如图所示：
	\onepicture[w=6cm]{figs/21_an.png}

	\item
	$b$，$c$；$S_{1}$，$E$（或$S_{2}$，$E$）

\end{enumerate}
}
%% memo:
\memoanswer{
本题原卷及材料中未提供详解，待补充。
}

\item
%% number: 22
%% paperName: 2013年山东理综第 22 题 16 分
%% typeId: 6
%% type: 计算题
%% chapter: 运动和力
%% point: 牛顿定律的应用
%% method: 极值
%% score: 16
%% degree: 650
%% duplicateId: 0
%% body:
如图所示，一质量$m=0.4\Ukg$的小物块，以$v_{0}=2\Ums$的初速度，在与斜面成某一夹角的拉力$F$作用下，沿斜面向上做匀加速运动，经$t=2\Us$的时间物块由$A$点运动到$B$点，$A$、$B$之间的距离$L=10\Um$。已知斜面倾角$\theta=30\Udu$，物块与斜面之间的动摩擦因数$\mu=\sqrt 3/3$。重力加速度$g$取$10\Umsq$。
\begin{enumerate}
	\item
	求物块加速度的大小及到达$B$点时速度的大小。
	\item
	拉力$F$与斜面的夹角多大时，拉力$F$最小？拉力$F$的最小值是多少？
\end{enumerate}

\onepicture[w=5cm, align=right]{figs/22.png}

%% answer:
\jdanswer{
\begin{enumerate}
	\item
	$a=3\Umsq$，$v=8\Ums$

	\item
	$\alpha=30\Udu$，$F_{\min}=\frac{13\sqrt 3}{5}\UN$

\end{enumerate}
}
%% memo:
\memoanswer{
（1）设物块加速度的大小为$a$，到达B点时速度的大小为$v$，由运动学公式得

$L=v_{0}t+\frac{1}{2}at^{2}$①，

$v=v_{0}+at$②，

联立①②得$a=3\Umsq$③，

$v=8\Ums$④。

（2）设物块所受支持力为$F_{N}$，所受摩擦力为$F_{f}$，拉力与斜面间的夹角为$\alpha$，受力分析如图所示，由牛顿第二定律得

\onepicture[w=5.5cm]{figs/22_an.png}

$F\cos\alpha-mg\sin\theta-F_{f}=ma$⑤，

$F\sin\alpha+F_{N}-mg\cos\theta=0$⑥，

又$F_{f}=\mu F_{N}$⑦，

联立⑤⑥⑦式得

$F=\frac{mg(\sin\theta+\mu\cos\theta)+ma}{\cos\alpha+\mu\sin\alpha}$⑧，

由数学知识得

$\cos\alpha+\frac{\sqrt 3}{3}\sin\alpha=\frac{2\sqrt 3}{3}\sin(60\Udu+\alpha)$⑨，

由⑧⑨式可知对应F最小的夹角为$\alpha=30\Udu$⑩，

联立③⑧⑩式，代入数据得F的最小值为$F_{\min}=\frac{13\sqrt 3}{5}\UN$\Circled{11}。
}

\item
%% number: 23
%% paperName: 2013年山东理综第 23 题 18 分
%% typeId: 6
%% type: 计算题
%% chapter: 磁场
%% point: 带电粒子在磁场中的运动
%% method: 无
%% score: 18
%% degree: 650
%% duplicateId: 0
%% body:
如图所示，在坐标系$xOy$的第一、第三象限内存在相同的匀强磁场，磁场方向垂直于$xOy$平面向里；第四象限内有沿$y$轴正方向的匀强电场，电场强度大小为$E$。一带电量为$+q$、质量为$m$的粒子，自$y$轴的P点沿$x$轴正方向射入第四象限，经$x$轴上的$Q$点进入第一象限，随即撤去电场，以后仅保留磁场。已知OP$=d$，OQ$=2d$，不计粒子重力。
\begin{enumerate}
	\item
	求粒子过Q点时速度的大小和方向。
	\item
	若磁感应强度的大小为一定值$B_{0}$，粒子将以垂直$y$轴的方向进入第二象限，求$B_{0}$。
	\item
	若磁感应强度的大小为另一确定值，经过一段时间后粒子将再次经过Q点，且速度与第一次过Q点时相同，求该粒子相邻两次经过Q点所用的时间。
\end{enumerate}

\onepicture[w=6cm, align=right]{figs/23.png}

%% answer:
\jdanswer{
\begin{enumerate}
	\item
	$2\sqrt{\frac{qEd}{m}}$，$45\Udu$

	\item
	$\sqrt{\frac{mE}{2qd}}$

	\item
	$(2+\pi)\sqrt{\frac{2md}{qE}}$

\end{enumerate}
}
%% memo:
\memoanswer{
（1）设粒子在电场中运动的时间为$t_{0}$，加速度的大小为$a$，粒子的初速度为$v_{0}$，过Q点时速度的大小为$v$，沿$y$轴方向分速度的大小为$v_{y}$，速度与$x$轴正方向间的夹角为$\theta$，由牛顿第二定律得

$qE=ma$①，

由运动学公式得

$d=\frac{1}{2}at_{0}^{2}$②，

$2d=v_{0}t_{0}$③，

$v_{y}=at_{0}$④，

$v=\sqrt{v_{0}^{2}+v_{y}^{2}}$⑤，

$\tan\theta=\frac{v_{y}}{v_{0}}$⑥，

联立①②③④⑤⑥式得

$v=2\sqrt{\frac{qEd}{m}}$⑦，

$\theta=45\Udu$⑧。

（2）设粒子做圆周运动的半径为$R_{1}$，粒子在第一象限的运动轨迹如图所示，$O_{1}$为圆心，由几何关系可知$\triangle O_{1}OQ$为等腰直角三角形，得

$R_{1}=2\sqrt{2}d$⑨，

由牛顿第二定律得

$qvB_{0}=m\frac{v^{2}}{R_{1}}$⑩，

\onepicture[w=5.5cm]{figs/23_an1.png}

联立⑦⑨⑩式得$B_{0}=\sqrt{\frac{mE}{2qd}}$\Circled{11}。

（3）设粒子做圆周运动的半径为$R_{2}$，由几何分析（粒子运动的轨迹如图所示，$O_{2}$、$O_{2}'$是粒子做圆周运动的圆心，Q、F、G、H是轨迹与两坐标轴的交点，连接$O_{2}$、$O_{2}'$，由几何关系知，$O_{2}FGO_{2}'$和$O_{2}QHO_{2}'$均为矩形，进而知FQ、GH均为直径，QFGH也是矩形，又$FH\perp GQ$，可知QFGH是正方形，$\triangle QOG$为等腰直角三角形）可知，粒子在第一、第三象限的轨迹均为半圆，得

$2R_{2}=2\sqrt{2}d$\Circled{12}，

粒子在第二、第四象限的轨迹为长度相等的线段，得

$FG=HQ=2R_{2}$\Circled{13}，

设粒子相邻两次经过Q点所用的时间为$t$，则有

$t=\frac{FG+HQ+2\pi R_{2}}{v}$\Circled{14}，

联立⑦\Circled{12}\Circled{13}\Circled{14}得

$t=(2+\pi)\sqrt{\frac{2md}{qE}}$\Circled{15}。

\onepicture[w=5.5cm]{figs/23_an2.png}
}

\item
%% number: 36
%% paperName: 2013年山东理综第 36 题 4 分
%% typeId: 6
%% type: 计算题
%% chapter: 气体性质
%% point: 理想气体状态方程
%% method: 无
%% score: 4
%% degree: 650
%% duplicateId: 0
%% body:
（1）下列关于热现象的描述正确的是\xzanswer{C}

\fourchoices[answer=C]
{根据热力学定律，热机的效率可以达到$100\%$}
{做功和热传递都是通过能量转化的方式改变系统内能的}
{温度是描述热运动的物理量，一个系统与另一个系统达到热平衡时两系统温度相同}
{物体由大量分子组成，其单个分子的运动是无规则的，大量分子的运动也是无规律的}

（2）我国“蛟龙”号深海探测船载人下潜超七千米，再创载人深潜新纪录。在某次深潜实验中，“蛟龙”号探测到$990\Um$深处的海水温度为$280\UK$。某同学利用该数据来研究气体状态随海水深度的变化，如图所示，导热良好的气缸内封闭一定质量的气体，不计活塞的质量和摩擦，气缸所处海平面的温度$T_{0}=300\UK$，压强$p_{0}=1\Uatm$，封闭气体的体积$V_{0}=3\Umc$。如果将该气缸下潜至$990\Um$深处，此过程中封闭气体可视为理想气体。
\begin{enumerate}
	\item
	求$990\Um$深处封闭气体的体积（$1\Uatm$相当于$10\Um$深的海水产生的压强）。
	\item
	下潜过程中封闭气体\tkanswer[1.2cm]{放热}（填“吸热”或“放热”），传递的热量\tkanswer[1.2cm]{大于}（填“大于”或“小于”）外界对气体所做的功。
\end{enumerate}

\onepicture[h=4cm, align=right]{figs/36.png}

%% answer:
\jdanswer{
\begin{enumerate}
	\item
	$V=2.8\times10^{-2}\Umc$

	\item
	放热，大于

\end{enumerate}
}
%% memo:
\memoanswer{
（2）①当气缸下潜至$990\Um$时，设封闭气体的压强为$p$，温度为$T$，体积为$V$，由题意知$p=100\Uatm$①，

根据理想气体状态方程得$\frac{p_{0}V_{0}}{T_{0}}=\frac{pV}{T}$②，

代入数据得$V=2.8\times10^{-2}\Umc$③。
}

\item
%% number: 37
%% paperName: 2013年山东理综第 37 题 4 分
%% typeId: 6
%% type: 计算题
%% chapter: 光学
%% point: 光的折射
%% method: 无
%% score: 4
%% degree: 650
%% duplicateId: 0
%% body:
（1）如图甲所示，在某一均匀介质中，A、B是振动情况完全相同的两个波源，其简谐运动表达式为$x=0.1\pi\sin(20\pi t)\Um$，介质中P点与A、B两波源间距离分别为$4\Um$和$5\Um$，两波源形成的简谐横波分别沿AP、BP方向传播，波速都是$10\Ums$。
\begin{enumerate}
	\item
	求简谐横波的波长。
	\item
	$P$点的振动\tkanswer[1.2cm]{加强}（填“加强”或“减弱”）。
\end{enumerate}

\onepicture[num=甲, w=4.5cm, align=right]{figs/37a.png}

（2）如图乙所示，ABCD是一直角梯形棱镜的横截面，位于截面所在平面内的一束光线由O点垂直AD边射入。已知棱镜的折射率$n=\sqrt 2$，AB$=$BC$=8\Ucm$，OA$=2\Ucm$，$\angle$OAB$=60\Udu$。
\begin{enumerate}
	\item
	求光线第一次射出棱镜时，出射光线的方向。
	\item
	第一次的出射点距C\tkanswer[1.2cm]{$\frac{4\sqrt 3}{3}$}$\Ucm$。
\end{enumerate}

\onepicture[num=乙, h=4cm, align=right]{figs/37b.png}

%% answer:
\jdanswer{
\begin{enumerate}
	\item
	$1\Um$

	\item
	加强

	\item
	从$CD$边射出，与$CD$边成$45\Udu$斜向左下方

	\item
	$\frac{4\sqrt 3}{3}$

\end{enumerate}
}
%% memo:
\memoanswer{
（1）解：设简谐波的波速为$v$，波长为$\lambda$，周期为$T$，由题意知

$T=0.1\Us$①，

由波速公式$v=\frac{\lambda}{T}$②，

代入数据得$\lambda=1\Um$③。

$P$点的振动加强。

（2）设发生全反射的临界角为$C$，由折射定律得

$\sin C=\frac{1}{n}$④，

代入数据得

$C=45\Udu$⑤，

光路图如图所示，由几何关系可知光线在AB边和BC边的入射角均为$60\Udu$，均发生全反射。设光线在CD边的入射角为$\alpha$，折射角为$\beta$，由几何关系得$\alpha=30\Udu$，小于临界角，光线第一次射出棱镜是在CD边，由折射定律得

$n=\frac{\sin\beta}{\sin\alpha}$⑥，

代入数据得$\beta=45\Udu$⑦。

\onepicture[w=5cm]{figs/37_an.png}
}

\item
%% number: 38
%% paperName: 2013年山东理综第 38 题 4 分
%% typeId: 6
%% type: 计算题
%% chapter: 运动和力
%% point: 动量守恒定律
%% method: 无
%% score: 4
%% degree: 650
%% duplicateId: 0
%% body:
（1）恒星向外辐射的能量来自于其内部发生的各种热核反应，当温度达到$10^{8}\UK$时，可以发生“氦燃烧”。完成“氦燃烧”的核反应方程：\ce{^{4}_{2}He} + \tkanswer[1.6cm]{$\ce{^{4}_{2}He}$或$\alpha$} $\to$ \ce{^{8}_{4}Be + \gamma}。\ce{^{8}_{4}Be}是一种不稳定的粒子，其半衰期为$2.6\times10^{-16}\Us$。一定质量的\ce{^{8}_{4}Be}，经$7.8\times10^{-16}\Us$后所剩\ce{^{8}_{4}Be}占开始时的\tkanswer[1.6cm]{$\frac{1}{8}$或$12.5\%$}。

（2）如图所示，光滑水平轨道上放置长板A（上表面粗糙）和滑块C，滑块B置于A的左端，三者质量分别为$m_{A}=2\Ukg$、$m_{B}=1\Ukg$、$m_{C}=2\Ukg$。开始时\textbf{$C$}静止，A、B一起以$v_{0}=5\Ums$的速度匀速向右运动，A与C发生碰撞（时间极短）后C向右运动，经过一段时间A、B再次达到共同速度一起向右运动，且恰好不再与C碰撞。求A与C发生碰撞后瞬间A的速度大小。

\onepicture[w=7cm, align=right]{figs/38.png}

%% answer:
\jdanswer{
（1）\ce{^{4}_{2}He}（或$\alpha$），$\frac{1}{8}$或$12.5\%$；（2）$v_{A}=2\Ums$
}
%% memo:
\memoanswer{
因碰撞时间极短，A与C碰撞过程动量守恒，设碰后瞬间A的速度为$v_{A}$，C的速度为$v_{C}$，以向右为正方向，由动量守恒定律得

$m_{A}v_{0}=m_{A}v_{A}+m_{C}v_{C}$①，

A与B在摩擦力作用下达到共同速度，设共同速度为$v_{AB}$，由动量守恒定律得

$m_{A}v_{A}+m_{B}v_{0}=(m_{A}+m_{B})v_{AB}$②，

A与B达到共同速度后恰好不再与C碰撞，应满足

$v_{AB}=v_{C}$③，

联立①②③式，代入数据得

$v_{A}=2\Ums$④。
}

\end{enumerate}

\end{document}
